\section{Impact and benefits of the project}

\subsection{Scientific impact}

\paragraph{Publications:}
Publication of results in the grade A peer-review journal 
Astronomy and Astrophysics (A\&A). 
We will ensure that open-access versions of our publications
will be published on the arXiv website, as required by the 
CNRS publication policy. 

\paragraph{Conferences:} 
Our results will be advertised in the main cosmological 
conferences around the world. Presentations are expected by 
the two new members of the group as well as the scientific
coordinator. These presentations will also contribute to 
the exposure of \projacro\ members as they search their 
next scientific opportunities. 

\paragraph{Organisation of a local CosmicFlows conference:}
During the project, we plan to organise a version of the
CosmicFlows conference locally in Marseille. This conference 
unites experts from the field of cosmology and astrophysics 
with peculiar velocities. The scientific coordinator was 
one of the invited speakers of the last version of the 
CosmicFlows\footnote{\url{https://indico.global/event/12931/}}
in 2025 in Brisbane, Australia, tying up good links with their 
organizers. 

\subsection{Societal impact}

\paragraph{Press-releases:}
Major results within the DESI or the ZTF collaborations 
can be part of a coordinated world-wide press-releases. 
Those press-releases are usually republished by main 
news papers such as Le Monde, Le Figaro, among others, 
significantly increasing the impact of \projacro\ to the
general public. 

\paragraph{Outreach events:}
Reaching out to the general public is one of the priorities of 
\projacro\. The members of the team will be invited to 
contribute to various outreach events in town and connect to 
the next generation of potential researchers. 

We will continue participating to local outreach events such as  
\emph{La Fête de la Science}, Pint of Science, as well as 
promoting our science through workshops for high-school students. 


